\section{Other useful Custom Exceptions}\label{sec:UsefulCustomExceptions}
Although the headline of this chapter is talking about \bdq{Custom Exceptions}, it deals more with custom implementations of \lstinline|java.lang.Error|\index{java.lang!Error} than those of \lstinline|java.lang.Exception|\index{java.lang!Exception}.

The error classes\index{java.lang!Error} \lstinline|UnlikelyExceptionError| and \lstinline|ImpossibleExceptionError| are mentioned in chapter \tqfullvref{sec:ImpossibleExceptions} where I stated that an \lstinline|UnlikelyExceptionError| should be thrown if an exception is caught that is theoretically possible at that particular point, but highly unlikely.

An \lstinline|ImpossibleExceptionError| should be thrown when an exception is caught that, in the given context, should not actually be able to occur.

Implementations for these exceptions are given in chapters \tqvref{code:ImpossibleExceptionError} and \tqvref{code:UnlikelyExceptionError}.

\bigskip

Chapter \tqfullvref{sec:TheSwitchStatement} stipulates that every \lstinline|switch|\index{switch} statement\index{switch!statement} (and, in principle, every \lstinline|switch|\index{switch} expression\index{switch!expression} as well) should include a \lstinline|default|\index{switch!default} branch, even if all possible cases are already covered by corresponding \lstinline|case|\index{switch!case} branches.

But what to put into that \lstinline|default|\index{switch!default} branch?

In case the \lstinline|switch|\index{switch} uses an \lstinline|enum| value as the selector, the recommendation is to throw an exception (more specific, an error) in the \lstinline|default|\index{switch!default} branch:
\keeplisting{1}
\begin{lstlisting}
switch( direction )
{
    case null -> throw new NullArgumentException( "direction" );
    case NORTH -> goUp();
    case SOUTH -> goDown();
    case WEST -> goLeft();
    case EAST -> goRight();
    default -> throw new UnsupportedEnumError( direction );
}
\end{lstlisting}

This protects your code from changes to the enum if that is not part of your project\footnote{Otherwise the compiler would complain: \bdq{\textit{'switch' expression does not cover all possible input values}}.}. Admittedly, it is rather unlikely that a fifth cardinal direction will ever be added; in that sense, the example is perhaps poorly chosen.

In chapter \tqfullvref{code:UnsupportedEnumError} you find an implementation of that error.
 
\bigskip

You may have noticed that the private constructor\index{constructor} for a utility class\index{class!utility class} that I gave as an example in chapter \tqvref{sec:ConstructorComment} throws an error, in particular, it throws an \lstinline|PrivateConstructorForStaticClassCalledError|.

Utility classes\index{class!utility class} are not meant to be instantiated, therefore they have only a single private constructor\index{constructor}. Nevertheless, that constructor could be called through reflection\index{reflection}, but when this constructor\index{constructor} does nothing else than throwing an error, it is now really impossible to instantiate that class.

The code for the \lstinline|PrivateConstructorForStaticClassCalledError| can be found in chapter \tqfullvref{code:PrivateConstructorForStaticClassCalledError}.

%==============================================================================
% End of file!!
%==============================================================================

